\section{Introduction}

The expansion of the Internet of Things (IoT), together with the deployment of data-intensive applications over 5G and emerging 6G networks, has increased the demand for low-latency, high-throughput, and reliable data access at the network edge. Applications such as smart city monitoring, autonomous transportation, industrial automation, augmented reality, and large-scale sensor networks continuously generate and consume delay-sensitive and context-dependent data. Depending on centralized cloud infrastructures to handle these workloads results in significant latency, backhaul congestion, and scalability bottlenecks.

Edge computing has emerged as a key architectural paradigm to address these issues by bringing computation and storage resources closer to end users. Among various edge computing mechanisms, \emph{content caching at the edge} plays a central role in reducing end-to-end latency and alleviating network load by storing frequently accessed data near consumers. Prior work has demonstrated that cooperative and distributed caching can significantly improve network performance in wireless systems~\cite{Shanmugam2013} and that proactive caching can further enhance quality of service under heavy and bursty traffic conditions~\cite{Bastug2014}.

However, caching in IoT-driven edge environments presents unique challenges that distinguish it from traditional web or content delivery network (CDN) caching~\cite{Podlipnig2003}. IoT systems are characterized by high user mobility, heterogeneous devices, non-stationary demand, and strict resource constraints. In such settings, data access patterns are often influenced by spatial locality, temporal dynamics, and user movement across edge nodes. As a result, caching decisions made independently at each node may lead to redundant content replication, inefficient storage utilization, and degraded system-wide performance.

To mitigate these limitations, prior research has explored cooperative caching mechanisms that allow neighboring edge nodes to coordinate content placement~\cite{Shanmugam2013,Poularakis2014}. While cooperative strategies provide clear improvements over purely local policies, most existing approaches assume static network topology or rely on predefined neighborhood structures. Even recent machine learning–based caching strategies~\cite{Shuja2021JNCA} typically focus on local demand prediction or reinforcement learning at individual nodes, treating topology implicitly or as a secondary feature rather than as a first-class object.

In parallel, graph-based learning methods particularly Graph Neural Networks (GNNs) have demonstrated strong capability in modeling relational dependencies in communication networks, including routing and resource allocation~\cite{Rusek2019,Jiang2022Survey}. These methods naturally encode structural information and generalize across varying network topologies. Despite this promise, the application of graph learning to cooperative caching in IoT edge environments, especially under mobility-induced topology dynamics, remains limited.

This work is motivated by the observation that user mobility implicitly defines a time-varying interaction graph among edge nodes. Nodes that frequently exchange users are likely to experience correlated demand patterns. Leveraging this mobility-derived topological structure through learning-based approaches may allow for more effective and coordinated caching decisions. Therefore, understanding how topological context influences cooperative caching behavior is therefore essential for designing scalable and adaptive edge systems in future IoT and 6G networks.